\documentclass[11pt]{article}
\usepackage[a4paper,margin=0.82in]{geometry}
\usepackage{microtype}
\usepackage[T1]{fontenc}
\usepackage{lmodern}
\usepackage{parskip}
\usepackage[hidelinks]{hyperref}
\pagestyle{empty}
\setlength{\emergencystretch}{2em}
\setlength{\parskip}{0.45em}

\begin{document}

{\raggedright
\textbf{Cristiano Capone}\\
\href{mailto:cristiano0capone@gmail.com}{cristiano0capone@gmail.com}\par}

\vspace{0.6em}
{\raggedright The Editors, \textit{Advanced Intelligent Systems}\par}

\vspace{0.8em}

Dear Editors,

We are pleased to submit our manuscript, ``Optimal stimulation sites are not the
most affected: personalised fMRI models of Alzheimer's disease,'' by Cristiano
Capone, Enza Cece, Luca Falorsi, Andrea Ciardiello, Guido Gigante, Evaristo
Cisbani and Maurizio Mattia, for consideration in \textit{Advanced Intelligent
Systems}.

Targeted neuromodulation rests on two decisions that are still made
heuristically: which site to engage, and at what amplitude. Both are typically
settled anatomically and fixed in advance. We ask whether they can instead be
\emph{computed}, for each individual, from a model of that person's own brain.

We build individualised digital twins: reservoir-computing models fit to
resting-state functional MRI so that their free-running dynamics reproduce each
individual's own lagged connectivity. The design choice that makes this work is
that the recurrent substrate is fixed and shared across every subject, and only a
linear read-out is fit per person. Three consequences follow. Per-subject
estimation stays tractable where a subject-specific autoregression would be
underdetermined, since $121$ regions must be fit from roughly $140$ volumes. All
fitted models occupy the same coordinates, so they can be compared, averaged and
interpolated. And the reservoir's separate input pathway makes ``drive region
$k$'' an explicit, well-posed operation rather than an ad hoc addition. The
resulting twins are identifiable: a model fit to one individual describes that
individual's held-out data better than anyone else's.

Perturbing them yields a result with a direct engineering consequence. Correcting
a twin's connectivity toward a control template reverts the disease signature,
but the necessary correction is distributed, and no single-node edit reproduces
it. A focal drive at the region whose connectivity is most altered fails at any
amplitude, whereas a drive at the region selected by its \emph{effect on the
disease discriminant} succeeds from one site. The two site sets are nearly
disjoint, and the reason is mechanistic rather than incidental: the most-affected
regions are the most weakly connected, and therefore poor actuators. Targets are
individual, and a real-time closed-loop controller that titrates amplitude from
the model's own output reaches comparable efficacy at substantially lower dose,
using only causally available information.

We are explicit about scope. This is an in-silico study; no patient was
stimulated, and complete reversal requires amplitudes above the physiological
range. The contribution is a computable, per-patient prescription for where and
how hard to drive, and a demonstration that the intuitive target is
systematically the wrong one. We believe this speaks to readers of
\textit{Advanced Intelligent Systems} working on digital twins, closed-loop
control and model-based decision-making in biological systems, beyond the
Alzheimer's application in which we demonstrate it.

The manuscript is original, is not under consideration elsewhere, and all authors
have approved its submission. The data derive from the Alzheimer's Disease
Neuroimaging Initiative and are used under its Data Use Agreement. Code is
available at
\href{https://github.com/cristianocapone/AD-reservoir-FC-stimulation}%
{github.com/cristianocapone/AD-reservoir-FC-stimulation}. We declare no
competing interests, and have suggested potential referees who have no conflict
of interest with the authors. We thank you for considering our work.

\vspace{1em}
{\raggedright
Sincerely,\\[0.3em]
Cristiano Capone, on behalf of all authors\par}

\end{document}
